\chapter{Tracheostomy Management}

\section*{What Is This Test or Procedure?}
Tracheostomy management is the ongoing care of a tracheostomy tube, including cleaning, suctioning, cuff management, tube changes, and prevention of complications. It is required for patients with a surgical airway for prolonged ventilation, airway obstruction, or secretion clearance.
\section*{Alternative Names}
Tracheostomy care; trach care; tracheostomy tube management; stoma care
\section*{Principle}
A tracheostomy tube bypasses the upper airway, delivering air directly to the trachea. Management maintains tube patency by removing secretions, secures the tube in place, provides humidification, and monitors the stoma and cuff pressures to prevent tissue injury and infection.
\section*{History}
Tracheostomy has been performed since antiquity, but modern indwelling tracheostomy tubes and structured management protocols developed during the 20th century, particularly with the growth of intensive care and prolonged mechanical ventilation.
\section*{Why This Test or Procedure Is Ordered}
Required for patients who cannot tolerate decannulation, those on long-term mechanical ventilation, patients with upper airway obstruction or aspiration, and as part of routine care to prevent tube occlusion and infection.
\section*{Is This a Primary or Secondary Test or Procedure}
A primary supportive procedure that maintains a patent airway and is essential to the safe ongoing use of a tracheostomy.
\section*{How This Test or Procedure Is Performed}
Routine care includes humidified oxygen or air, regular suctioning with a sterile catheter, cleaning of the inner cannula, and skin care around the stoma. The tube ties are changed and the tube is secured, the cuff is deflated periodically and its pressure checked, and tube changes are performed by trained staff when needed.
\section*{What Preparation Is Needed}
Gathering suction, cleaning supplies, a replacement tube of the same size, ties, and saline before any procedure. The patient is positioned with the head slightly extended, and the assistant is prepared for emergency tube reinsertion.
\section*{Contraindications and Precautions}
Avoid disconnecting humidification for prolonged periods, and never advance suction beyond the pre-measured length. Precautions include confirming a mature stoma before tube changes, keeping a spare tube and obturator at the bedside, and having an emergency plan for accidental decannulation.
\section*{Risks}
Tube occlusion, accidental decannulation, bleeding, infection, tracheal stenosis, and tracheal wall injury from high cuff pressures.
\section*{What Happens After the Test or Procedure}
The patient's oxygen saturation and breathing are monitored after suctioning or tube changes, and the tube position is confirmed. Any change in ventilation, bleeding, or signs of infection is reported promptly.
\section*{Understanding Your Results}
A patent airway with clear breath sounds, stable oxygenation, and a clean stoma indicates effective management. Increasing suction requirements, fever, or bleeding may indicate a complication.
\section*{Normal Results}
A clean, dry, non-inflamed stoma with a secure tracheostomy tube, clear airways, and stable respiratory status.
\section*{Follow-up Tests and Procedures}
The tracheostomy is reviewed by the clinical team, and the patient is assessed for readiness for decannulation with a trial of capping when the underlying indication has resolved. Speech and swallowing assessment and bronchoscopy may be needed.
\section*{References}
\begin{enumerate}
  \item MedlinePlus. Tracheostomy - aftercare. U.S. National Library of Medicine; 2025.
  \item Current Medical Diagnosis and Treatment. McGraw-Hill; 2025.
\end{enumerate}

\clearpage
